\documentclass[11pt,a4paper]{article}

% ============================================================================
% PACKAGES
% ============================================================================

\usepackage[utf8]{inputenc}
\usepackage[T1]{fontenc}
\usepackage{amsmath,amssymb,amsthm}
\usepackage{graphicx}
\usepackage{hyperref}
\usepackage{url}
\usepackage{cite}
\usepackage{algorithm}
\usepackage{algorithmic}
\usepackage{listings}
\usepackage{xcolor}
\usepackage{geometry}
\usepackage{titlesec}
\usepackage{enumitem}
\usepackage{booktabs}
\usepackage{multirow}
\usepackage{array}

% ============================================================================
% PAGE SETUP
% ============================================================================

\geometry{
    left=2.5cm,
    right=2.5cm,
    top=3cm,
    bottom=3cm
}

% ============================================================================
% HYPERREF SETUP
% ============================================================================

\hypersetup{
    colorlinks=true,
    linkcolor=blue,
    filecolor=magenta,
    urlcolor=cyan,
    citecolor=blue,
    pdftitle={Mechanism Design for Sustainable Hook Marketplace with Public Bytecode},
    pdfauthor={Hook Bazaar Research Team},
    pdfsubject={Mechanism Design, Game Theory, Blockchain, Smart Contracts}
}

% ============================================================================
% THEOREM ENVIRONMENTS
% ============================================================================

\theoremstyle{definition}
\newtheorem{definition}{Definition}[section]
\newtheorem{theorem}{Theorem}[section]
\newtheorem{lemma}{Lemma}[section]
\newtheorem{proposition}{Proposition}[section]
\newtheorem{corollary}{Corollary}[section]
\newtheorem{example}{Example}[section]

\theoremstyle{remark}
\newtheorem{remark}{Remark}[section]

% ============================================================================
% CUSTOM COMMANDS
% ============================================================================

% Game theory notation
\newcommand{\payoff}[2]{u_{#1}(#2)}
\newcommand{\strategy}[1]{s_{#1}}
\newcommand{\equilibrium}{\text{NE}}
\newcommand{\nash}{\text{Nash}}

% Mechanism design notation
\newcommand{\mechanism}{\mathcal{M}}
\newcommand{\type}{\theta}
\newcommand{\allocation}{a}
\newcommand{\payment}{p}

% Hook marketplace notation
\newcommand{\hook}{\text{hook}}
\newcommand{\license}{\text{license}}
\newcommand{\developer}{\text{dev}}
\newcommand{\buyer}{\text{buyer}}
\newcommand{\marketplace}{\text{marketplace}}

% ============================================================================
% TITLE AND AUTHORS
% ============================================================================

\title{%
    \textbf{Mechanism Design for Sustainable Hook Marketplace} \\
    \large{Protecting Intellectual Property with Public Bytecode}
}

\author{%
    Hook Bazaar Research Team\thanks{Corresponding author: research@hookbazaar.xyz} \\
    \textit{Hook Bazaar Protocol}
}

\date{\today}

% ============================================================================
% DOCUMENT
% ============================================================================

\begin{document}

\maketitle

% ============================================================================
% ABSTRACT
% ============================================================================

\begin{abstract}
This paper presents a mechanism design framework for creating a sustainable marketplace for Uniswap v4 hooks where bytecode is publicly visible on-chain, yet purchasing licenses remains economically superior to decompiling code. We address the fundamental challenge that if decompilation costs less than purchasing a license, the marketplace becomes unsustainable. 

Our approach combines game-theoretic mechanisms, economic incentives, and cryptographic techniques to create conditions where honest participation is the dominant strategy. We propose a dual-deposit escrow mechanism for license transactions, a reputation system based on PeerTrust for building trust, and integration with Fhenix confidential computing for protecting sensitive parameters. Additionally, we design an auditor incentive mechanism to prevent information leakage and implement bytecode encryption to make decompilation economically unviable.

We prove that under our proposed mechanisms, purchasing licenses becomes the subgame perfect Nash equilibrium, ensuring market sustainability. Our framework provides both theoretical guarantees and practical implementation strategies, making it suitable for real-world deployment in blockchain-based software marketplaces.
\end{abstract}

\textbf{Keywords:} Mechanism Design, Game Theory, Blockchain, Smart Contracts, Intellectual Property, Zero-Knowledge Proofs, Confidential Computing

% ============================================================================
% TABLE OF CONTENTS
% ============================================================================

\newpage
\tableofcontents
\newpage

% ============================================================================
% INTRODUCTION
% ============================================================================

\section{Introduction}

\subsection{Motivation}

The emergence of programmable blockchains has enabled new forms of software marketplaces where code is deployed on-chain and publicly visible. This transparency creates a fundamental challenge for intellectual property protection: if bytecode is publicly available, what prevents users from decompiling and using code without purchasing licenses?

Hook Bazaar is a marketplace for Uniswap v4 hooks that enables developers to monetize their implementations while allowing protocol designers to purchase licenses for specific pools. The marketplace must address the following core dilemma:

\begin{enumerate}
    \item \textbf{Transparency Requirement}: Bytecode must be deployed on-chain for verification and execution
    \item \textbf{Protection Requirement}: Developers must be able to monetize their intellectual property
    \item \textbf{Sustainability Requirement}: Purchasing licenses must be economically superior to decompilation
\end{enumerate}

\subsection{Problem Statement}

Formally, we define the sustainability condition:

\begin{definition}[Market Sustainability Condition]
A hook marketplace is sustainable if and only if:
\begin{equation}
P_{\text{purchase}} < C_{\text{decompile}} + C_{\text{opportunity}} + R_{\text{risk}} + V_{\text{added}}
\end{equation}
where:
\begin{itemize}
    \item $P_{\text{purchase}}$ is the license purchase price
    \item $C_{\text{decompile}}$ is the cost of decompiling bytecode
    \item $C_{\text{opportunity}}$ is the opportunity cost of decompilation time
    \item $R_{\text{risk}}$ is the risk premium (legal, security, etc.)
    \item $V_{\text{added}}$ is the value of added services (documentation, audits, support)
\end{itemize}
\end{definition}

\subsection{Contributions}

This paper makes the following contributions:

\begin{enumerate}
    \item \textbf{Theoretical Framework}: We provide a game-theoretic analysis of hook marketplace sustainability with public bytecode, proving conditions for Nash equilibrium where purchasing is the dominant strategy.
    
    \item \textbf{Economic Mechanisms}: We design multiple economic mechanisms including dual-deposit escrow, reputation systems, and value-added services that collectively make purchasing superior to decompilation.
    
    \item \textbf{Cryptographic Solutions}: We integrate Fhenix confidential computing for parameter protection and bytecode encryption, making decompilation cryptographically infeasible.
    
    \item \textbf{Auditor Incentive Design}: We propose a mechanism to prevent auditor information leakage using dual deposits and cryptographic watermarking.
    
    \item \textbf{Practical Implementation}: We provide detailed specifications for smart contract implementation and integration strategies.
\end{enumerate}

\subsection{Paper Organization}

The remainder of this paper is organized as follows:
\begin{itemize}
    \item Section~\ref{sec:related} reviews related work in mechanism design, game theory, and blockchain marketplaces
    \item Section~\ref{sec:model} presents our formal model of the hook marketplace
    \item Section~\ref{sec:mechanisms} describes the economic mechanisms we propose
    \item Section~\ref{sec:analysis} provides game-theoretic analysis and proofs
    \item Section~\ref{sec:implementation} details implementation strategies
    \item Section~\ref{sec:security} analyzes security properties and attack vectors
    \item Section~\ref{sec:conclusion} concludes and discusses future work
\end{itemize}

% ============================================================================
% RELATED WORK
% ============================================================================

\section{Related Work}
\label{sec:related}

\subsection{Mechanism Design for Digital Goods}

The problem of selling digital goods without trusted third parties has been studied extensively. Asgaonkar and Krishnamachari~\cite{asgaonkar2018} propose a dual-deposit escrow mechanism for cheat-proof delivery of digital goods. We adapt their approach for hook licensing, extending it to handle the specific challenges of on-chain bytecode visibility.

\subsection{Reputation Systems in Marketplaces}

Reputation systems are crucial for reducing information asymmetry in marketplaces. Yamamoto and Hayashi~\cite{yamamoto2025} demonstrate that PeerTrust achieves the strongest alignment between price and quality in data trading markets. We apply similar principles to hook marketplaces, adapting the reputation system for software licensing contexts.

\subsection{Blockchain-Based Software Licensing}

Previous work on blockchain software licensing has focused primarily on token-based models or NFT licenses. Our approach differs by treating hooks as functional software components requiring ongoing support and verification, rather than as collectible assets.

\subsection{Confidential Computing for IP Protection}

Trusted Execution Environments (TEEs) and confidential computing platforms like Fhenix provide mechanisms for executing code in isolated environments. We leverage these technologies to protect sensitive parameters while maintaining public bytecode for verification.

\subsection{Zero-Knowledge Proofs for Code Verification}

Recent work has explored using zero-knowledge proofs to verify code correctness without revealing implementation~\cite{raziel2018}. We consider ZK proofs as a complementary approach to Fhenix for maximum protection.

% ============================================================================
% FORMAL MODEL
% ============================================================================

\section{Formal Model}
\label{sec:model}

\subsection{Players and Strategies}

We model the hook marketplace as a sequential game with three types of players:

\begin{definition}[Players]
\begin{itemize}
    \item \textbf{Developers} ($D$): Create and list hooks, set prices, invest in reputation
    \item \textbf{Buyers} ($B$): Purchase licenses or attempt decompilation
    \item \textbf{Marketplace} ($M$): Enforces mechanisms, maintains reputation system
\end{itemize}
\end{definition}

\subsection{Developer Utility Function}

The developer's utility is:
\begin{equation}
U_{\developer}(P, Q, R) = P - C_{\developer}(Q) - C_{\text{reputation}}(R) + \alpha \cdot R \cdot V_{\text{network}}(P, Q) - \text{Risk}(P, Q)
\end{equation}

where:
\begin{itemize}
    \item $P$ is the license price
    \item $Q$ is the hook quality
    \item $R$ is the developer reputation
    \item $\alpha$ is the network effect coefficient
    \item $V_{\text{network}}$ is the network value
\end{itemize}

\subsection{Buyer Utility Function}

The buyer's utility for purchasing is:
\begin{equation}
U_{\buyer}^{\text{purchase}}(P, \text{Trust}, Q) = V_{\hook}(Q) - P - C_{\text{integration}} + \beta \cdot \text{Trust} \cdot V_{\text{network}} + V_{\text{added}} - R_{\text{legal}}
\end{equation}

The buyer's utility for decompiling is:
\begin{equation}
U_{\buyer}^{\text{decompile}}(Q) = V_{\hook}(Q) - C_{\text{decompile}} - C_{\text{validation}} - C_{\text{legal\_risk}} - \gamma \cdot C_{\text{opportunity}}
\end{equation}

\subsection{Market Equilibrium}

\begin{definition}[Market Equilibrium]
A market equilibrium is a tuple $(P^*, Q^*, R^*, \text{Trust}^*)$ such that:
\begin{enumerate}
    \item $(P^*, Q^*, R^*)$ maximizes developer utility
    \item $U_{\buyer}^{\text{purchase}}(P^*, \text{Trust}^*, Q^*) > U_{\buyer}^{\text{decompile}}(Q^*)$
    \item $\text{Trust}^* = f(R^*, \text{MarketParticipation})$
\end{enumerate}
\end{definition}

% ============================================================================
% ECONOMIC MECHANISMS
% ============================================================================

\section{Economic Mechanisms}
\label{sec:mechanisms}

\subsection{Dual-Deposit Escrow}

We adapt the dual-deposit escrow mechanism from Asgaonkar and Krishnamachari~\cite{asgaonkar2018} for hook licensing.

\begin{theorem}[Subgame Perfect Nash Equilibrium]
In the dual-deposit escrow mechanism for hook licensing, honest behavior by both developer and buyer is the unique subgame perfect Nash equilibrium, provided that deposits $D_{\developer}$ and $D_{\buyer}$ are sufficiently large.
\end{theorem}

\begin{proof}
By backward induction, we show that:
\begin{enumerate}
    \item If buyer accepts delivery, both parties receive positive payoffs
    \item If buyer disputes, the smart contract verifies and penalizes the cheating party
    \item Given sufficiently large deposits, the expected payoff from cheating is negative
    \item Therefore, honest behavior is the dominant strategy
\end{enumerate}
\end{proof}

\subsection{Reputation System}

We implement a PeerTrust-based reputation system adapted for hooks:

\begin{equation}
R_i = \frac{\sum_{j} w_{ij} \cdot r_{ij}}{\sum_{j} w_{ij}}
\end{equation}

where $w_{ij}$ incorporates:
\begin{itemize}
    \item Time decay (recent reviews weighted more)
    \item Reviewer credibility
    \item Transaction context
    \item Review consistency
\end{itemize}

\subsection{Value-Added Services}

Licensed hooks include:
\begin{itemize}
    \item Comprehensive documentation (IPFS)
    \item Security audit reports
    \item Test suites and verification proofs
    \item Developer support and updates
    \item Legal protection and warranty
\end{itemize}

These services create value $V_{\text{added}}$ that decompiled code lacks.

\subsection{Network Effects}

Licensed hooks provide access to:
\begin{itemize}
    \item Marketplace ecosystem
    \item Compatibility guarantees
    \item Community support
    \item Partnership opportunities
\end{itemize}

The network value $V_{\text{network}}$ grows with market participation, creating positive feedback loops.

% ============================================================================
% GAME-THEORETIC ANALYSIS
% ============================================================================

\section{Game-Theoretic Analysis}
\label{sec:analysis}

\subsection{Nash Equilibrium Conditions}

\begin{theorem}[Sustainable Market Equilibrium]
A sustainable market equilibrium exists if and only if:
\begin{equation}
\beta \cdot \text{Trust} + C_{\text{validation}} + C_{\text{legal\_risk}} + \gamma \cdot C_{\text{opportunity}} > P - C_{\text{integration}} - C_{\text{decompile}}
\end{equation}
\end{theorem}

\begin{proof}
For buyers to prefer purchasing over decompiling:
\begin{align}
U_{\buyer}^{\text{purchase}} &> U_{\buyer}^{\text{decompile}} \\
V_{\hook} - P - C_{\text{integration}} + \beta \cdot \text{Trust} \cdot V_{\text{network}} + V_{\text{added}} - R_{\text{legal}} &> \\
&\quad V_{\hook} - C_{\text{decompile}} - C_{\text{validation}} - C_{\text{legal\_risk}} - \gamma \cdot C_{\text{opportunity}}
\end{align}

Rearranging and substituting $V_{\text{added}} = V_{\text{docs}} + V_{\text{audit}} + V_{\text{support}} + V_{\text{legal}}$:
\begin{equation}
\beta \cdot \text{Trust} \cdot V_{\text{network}} + V_{\text{added}} + C_{\text{validation}} + C_{\text{legal\_risk}} + \gamma \cdot C_{\text{opportunity}} > P - C_{\text{integration}} + C_{\text{decompile}}
\end{equation}

Under reasonable assumptions about value-added services and network effects, this inequality holds.
\end{proof}

\subsection{Mechanism Design Properties}

\begin{theorem}[Incentive Compatibility]
The proposed mechanism is incentive-compatible: truth-telling is a dominant strategy for all players.
\end{theorem}

\begin{theorem}[Individual Rationality]
All players receive non-negative utility from participation.
\end{theorem}

\begin{theorem}[Budget Balance]
The mechanism is budget-balanced: no external subsidy is required.
\end{theorem}

% ============================================================================
% CRYPTOGRAPHIC SOLUTIONS
% ============================================================================

\section{Cryptographic Solutions}
\label{sec:crypto}

\subsection{Fhenix Confidential Computing}

Fhenix enables execution of hook logic in encrypted environments where:
\begin{itemize}
    \item Sensitive parameters are encrypted
    \item Execution is isolated
    \item Results are encrypted
    \item Only licensed users can decrypt
\end{itemize}

\subsection{Bytecode Encryption}

We propose encrypting hook bytecode before deployment:
\begin{itemize}
    \item Bytecode encrypted using Fhenix homomorphic encryption
    \item Deployed as ciphertext on-chain
    \item Decryption keys granted only to licensed users
    \item Decompilation yields encrypted data, not source code
\end{itemize}

\subsection{Auditor Incentive Mechanism}

To prevent auditor information leakage:
\begin{itemize}
    \item Dual-deposit escrow for audits
    \item Cryptographic watermarking per auditor
    \item Leakage detection with slashing
    \item Reputation system for honest behavior
\end{itemize}

% ============================================================================
% IMPLEMENTATION
% ============================================================================

\section{Implementation}
\label{sec:implementation}

\subsection{Smart Contract Architecture}

The marketplace consists of:
\begin{itemize}
    \item \texttt{HookMarketplace}: Main marketplace contract
    \item \texttt{LicenseManager}: License lifecycle management
    \item \texttt{RevenueShare}: Revenue distribution
    \item \texttt{AuditorEscrow}: Auditor incentive mechanism
    \item \texttt{FhenixHookExecutor}: Confidential execution wrapper
\end{itemize}

\subsection{IPFS Metadata Structure}

Hook metadata stored on IPFS includes:
\begin{itemize}
    \item Documentation and specifications
    \item Security audit reports
    \item Test suites and verification proofs
    \item Developer reputation data
\end{itemize}

\subsection{Integration with Existing Architecture}

The mechanism integrates with:
\begin{itemize}
    \item \texttt{ProtocolAdminPanel}: Protocol management
    \item \texttt{MasterHook}: Hook execution
    \item \texttt{HookOrchestrator}: Mediator component
\end{itemize}

% ============================================================================
% SECURITY ANALYSIS
% ============================================================================

\section{Security Analysis}
\label{sec:security}

\subsection{Attack Vectors}

We analyze protection against:
\begin{itemize}
    \item Decompilation attacks
    \item License violations
    \item Auditor information leakage
    \item Reputation manipulation
    \item Parameter extraction
\end{itemize}

\subsection{Security Guarantees}

Our mechanisms provide:
\begin{itemize}
    \item Cryptographic protection (Fhenix encryption)
    \item Economic disincentives (deposits, slashing)
    \item Reputation-based trust
    \item Legal framework support
\end{itemize}

% ============================================================================
% CONCLUSION
% ============================================================================

\section{Conclusion}
\label{sec:conclusion}

This paper presents a comprehensive mechanism design framework for creating a sustainable hook marketplace despite public bytecode visibility. By combining game-theoretic mechanisms, economic incentives, and cryptographic techniques, we ensure that purchasing licenses is the dominant strategy, making the marketplace economically viable.

Our contributions include:
\begin{itemize}
    \item Theoretical framework with formal proofs
    \item Practical implementation specifications
    \item Security analysis and attack mitigation
    \item Integration strategies for real-world deployment
\end{itemize}

Future work includes:
\begin{itemize}
    \item Empirical evaluation of mechanism parameters
    \item Integration with zero-knowledge proofs
    \item Optimization of gas costs and performance
    \item Extension to other blockchain software marketplaces
\end{itemize}

% ============================================================================
% ACKNOWLEDGMENTS
% ============================================================================

\section*{Acknowledgments}

We thank the Hook Bazaar community for valuable feedback and discussions. This work is supported by the Hook Bazaar Protocol development team.

% ============================================================================
% BIBLIOGRAPHY
% ============================================================================

\begin{thebibliography}{99}

\bibitem{asgaonkar2018}
A. Asgaonkar and B. Krishnamachari.
\newblock Solving the Buyer and Seller's Dilemma: A Dual-Deposit Escrow Smart Contract for Provably Cheat-Proof Delivery and Payment for a Digital Good without a Trusted Mediator.
\newblock \textit{arXiv preprint arXiv:1806.08379}, 2018.

\bibitem{yamamoto2025}
K. Yamamoto and T. Hayashi.
\newblock Designing Reputation Systems for Manufacturing Data Trading Markets: A Multi-Agent Evaluation with Q-Learning and IRL-Estimated Utilities.
\newblock \textit{arXiv preprint arXiv:2511.19930}, 2025.

\bibitem{raziel2018}
D. Cerezo Sánchez.
\newblock Raziel: Private and Verifiable Smart Contracts on Blockchains.
\newblock \textit{arXiv preprint arXiv:1807.09484}, 2018.

\bibitem{akerlof1970}
G. A. Akerlof.
\newblock The Market for ``Lemons'': Quality Uncertainty and the Market Mechanism.
\newblock \textit{Quarterly Journal of Economics}, 84(3):488--500, 1970.

\bibitem{fudenberg1991}
D. Fudenberg and J. Tirole.
\newblock \textit{Game Theory}.
\newblock MIT Press, 1991.

\bibitem{xiong2004}
L. Xiong and L. Liu.
\newblock PeerTrust: Supporting Reputation-Based Trust for Peer-to-Peer Electronic Communities.
\newblock \textit{IEEE Transactions on Knowledge and Data Engineering}, 16(7):843--857, 2004.

\end{thebibliography}

% ============================================================================
% APPENDICES
% ============================================================================

\appendix

\section{Proof Details}
\label{app:proofs}

% Additional proof details can go here

\section{Implementation Code}
\label{app:code}

% Code listings can go here

\section{Additional Figures}
\label{app:figures}

% Additional figures can go here

\end{document}

